\solchapter{complex}

\sol{complex}{1} $1-1/R_0$: $33\%$, $67\%$, and $93\%$. A vaccine that protects $90\%$ of recipients leaves at most $90\%$ of the population immune even if everyone is vaccinated --- short of the $93\%$ measles needs. In practice a second dose raises protection to about $97\%$, which is why the measles vaccine is given twice.

\sol{complex}{2} $H_{100}/H_{50\,000} = (\ln100 + 0.577)/(\ln50\,000 + 0.577) = 5.19/11.4\approx45\%$. The hundred commonest words make up nearly half of everything you read.

\sol{complex}{3} $\int_1^\infty x\cdot x^{-\alpha}\dd x$ converges only for $\alpha>2$, and $\int_1^\infty x^2\cdot x^{-\alpha}\dd x$ only for $\alpha>3$. With infinite variance, the sample variance never settles: it keeps growing as you collect data, dominated by the largest event seen so far. Averages fluctuate wildly, the central limit theorem fails, and events far beyond anything in the historical record are to be expected.

\sol{complex}{4} $\dd s/\dd r = -\beta si/\gamma i = -R_0s$, so $s = e^{-R_0r}$ (starting from $s\approx1$, $r\approx0$). At the end, $i = 0$ and $s_\infty = 1 - r_\infty$, giving $r_\infty = 1 - e^{-R_0r_\infty}$. For $R_0 = 2$, iterating gives $r_\infty = 0.80$: $20\%$ are never infected. Note the overshoot: the epidemic stops growing once $50\%$ have been infected, but cases already in progress carry it to $80\%$.

\sol{complex}{5} A node is outside the giant component if each of its neighbours is, independently, outside it (through the rest of the network). Averaging $u^k$ over a Poisson degree distribution with mean $c$ gives $u = \sum_ke^{-c}c^ku^k/k! = e^{-c(1-u)}$, so $S = 1-u$ obeys $S = 1 - e^{-cS}$. The right side has slope $c$ at $S = 0$ and is concave, so it crosses the line $y = S$ at a positive value only if $c>1$. This is exactly the SIR final-size equation with $c$ in place of $R_0$: an epidemic is percolation on the network of transmissions.

\sol{complex}{6} Following a random edge reaches a node with probability proportional to its degree, so the expected degree found there is $\sum_kk\cdot kP(k)/\sum_kkP(k) = \langle k^2\rangle/\langle k\rangle$. Since $\langle k^2\rangle - \langle k\rangle^2 = \mathrm{Var}(k)\ge0$, this is at least $\langle k\rangle$, with equality only if every node has the same degree.

\sol{complex}{7} The energy released by quakes near magnitude $M$ goes as (number)$\times$(energy each) $\propto10^{-M}\cdot10^{1.5M} = 10^{0.5M}$, increasing with $M$. The total is dominated by the largest earthquakes: a single magnitude~9 releases more energy than all the smaller earthquakes of several decades combined.

\sol{complex}{8} Die: $\log_26 = 2.58$ bits. Coin: $1$ bit. Biased coin: $-0.9\log_20.9 - 0.1\log_20.1 = 0.47$ bits. Compressing to $25\%$ of $8$ bits shows English has at most $2$ bits per character; better models, and Shannon's own guessing experiments, suggest about $1$.

\sol{complex}{9} $\dd k_i/k_i = \dd t/2t$ gives $k_i(t) = m(t/t_i)^{1/2}$. A node has degree greater than $k$ if $t_i<m^2t/k^2$; since arrival times are uniform on $[0,t]$, the fraction is $m^2/k^2$. Differentiating, $P(k) = 2m^2/k^3$. The oldest nodes have had the longest to accumulate links, and since growth is proportional to degree, their head start compounds: rich get richer.

\sol{complex}{10} Random pairwise exchanges that conserve total money are mathematically like elastic collisions that conserve total energy. The equilibrium is the most probable distribution of a fixed total among many agents, which maximizes entropy subject to the constraint and is exponential, $P(m) = e^{-m/T}/T$, with ``temperature'' equal to the average, $T = \$100$. A simulation with $2000$ agents and $400\,000$ exchanges gives $64\%$ of agents below \$100 (exponential: $63\%$), $5.1\%$ above \$300 (exponential: $5.0\%$), and the richest $1\%$ holding $5.5\%$ of the money. In real economies the richest $1\%$ hold a third or more of the wealth: the top follows a Pareto power law, produced by multiplicative effects --- returns on capital proportional to wealth --- that a pure exchange model lacks.

\sol{complex}{11} A run on a $50\times50$ grid, discarding the first $15\,000$ grains and recording $60\,000$ more, gave a mean height of $2.10$ grains per site (the known value for large grids is $2.125$); $43\%$ of added grains triggered an avalanche, of sizes up to about $8400$ topplings. On log-log axes, the size distribution is a straight line over three decades, with slope about $-1.0$ to $-1.1$ (large-system studies give an exponent near $1.1$--$1.2$), cut off at sizes set by the grid. Why critical: if the pile is subcritical (too flat), grains pile up with few avalanches and the slope rises; if it is supercritical (too steep), large avalanches carry grains off the edges and the slope falls. Slow driving and fast relaxation pin the system to the boundary between the two, where avalanches of every size occur.
